\begin{abstract}
Ausgangspunkt ist eine elementare Umdeutung des Einheitskreises: Nicht Sinus
und Kosinus werden als primär behandelt, sondern ein linear fallender und
anschließend gespiegelter Kontrollwert. Durch einen algebraischen
Wurzellift erzeugt dieser Kontrollwert dennoch exakt einen Kreis oder eine
Ellipse. Das führt zu einer Trennung zwischen geometrischer Spur und zeitlicher
Parametrisierung. Eine gewichtete Reihe ungerader Kosinusharmonischer verbindet
die nichtglatte Dreieckswelle stetig mit der klassischen harmonischen
Parametrisierung. Der algebraische Invarianzkern und die atomare
Dreieckssteuerung werden in Lean~4 formalisiert. Die weiterführenden Aussagen
über Konvergenz, Glattheit und Grenzwerte bilden klar abgegrenzte analytische
Beweisverpflichtungen.
\end{abstract}
